% arara: xelatex: { shell: no }
% arara: bibtex
% arara: xelatex: { shell: no }
% arara: xelatex: { shell: no }
\documentclass[a4paper,11pt]{article}

% ---- Packages ----
\usepackage[utf8]{inputenc}
\usepackage[T1]{fontenc}
\usepackage{lmodern}
\usepackage{amsmath,amssymb,amsthm,bm}
\usepackage{mathtools}
\usepackage{booktabs}
\usepackage{graphicx}
\usepackage{siunitx}
\usepackage{hyperref}
\usepackage[margin=2.5cm]{geometry}
\usepackage{algpseudocode}
\usepackage{algorithm}
\usepackage{xcolor}
\usepackage{natbib}
\usepackage{caption}
\usepackage{subcaption}
\usepackage{enumitem}

% ---- Custom commands ----
\newcommand{\bq}{\,\text{Bq}}
\newcommand{\bqmt}{\,\text{Bq/m}^3}
\newcommand{\bqms}{\,\text{Bq/m}^2}
\newcommand{\keV}{\,\text{keV}}
\newcommand{\Eone}{\mathrm{E}_1}
\newcommand{\R}{\mathbb{R}}
\DeclareMathOperator*{\argmin}{arg\,min}
\DeclareMathOperator{\tr}{tr}

% ---- Title ----
\title{
  Geophysics-Inspired Fredholm Inversion for In-Situ Gamma\ Spectrometry:\
  Reconstructing Radionuclide Depth Profiles with Transport-Chemistry Priors
}

\author{Anonymous}
\date{August 2026}

\begin{document}

\maketitle

% =================================================================
\begin{abstract}
Determining the vertical distribution of radionuclides in soil from in-situ
gamma-spectrometry measurements requires solving a Fredholm integral
equation of the first kind --- a classically ill-posed problem.  This paper
presents the \texttt{soilactivity} software framework, which adapts modern
geophysical inversion methods to the specific structure of the gamma
transport kernel.  We introduce a comprehensive solver suite comprising
ten inversion algorithms: classical Tikhonov regularisation with
non-negativity constraints, truncated SVD with depth normalisation,
Landweber and Cimmino iterative methods, conjugate gradient least squares
(CGLS) with early stopping, randomised Kaczmarz, FISTA for
\(\ell_1\)-sparse inversion, total variation (TV) regularisation via ADMM,
and IRLS-based focusing inversion in both minimum-support and
minimum-gradient-support variants.  The parameter selection toolkit
includes six criteria: generalised cross-validation (GCV), L-curve corner
detection via Menger curvature, the Morozov discrepancy principle,
quasi-optimality, normalised cumulative periodogram (NCP), and a
signal-to-noise ratio (SNR) criterion.  Bayesian uncertainty quantification
is provided through ensemble Kalman inversion (EKI) and Laplace MAP
approximation.

A distinguishing feature of the framework is the integration of reactive
transport chemistry as an informed prior.  The transport module implements
the advection--dispersion equation (ADE) with linear, pH-dependent, and
Freundlich isotherm sorption models, competitive ion-exchange for
Cs-137 and Sr-90, Eh-dependent uranium speciation, two-site kinetic
sorption, Bateman decay chains, multi-layer soil columns, and a
Crank--Nicolson numerical solver.  Depth weighting strategies adapted
from gravity and magnetic inversion (Li--Oldenburg column normalisation,
power-law and logarithmic scaling, adaptive data-driven weighting) counter
the exponential sensitivity decay of the gamma kernel.  Joint multi-nuclide
inversion with coupling constraints and an automated ensemble pipeline
with AIC/BIC model selection complete the framework.

Twin experiments with synthetic Eu-152 and Cs-137 profiles demonstrate
that the geophysical solvers (TV, focusing MGS, CGLS) recover piecewise-
constant and layered contamination profiles substantially better than
classical smooth Tikhonov inversion, while the parametric transport model
provides physically interpretable parameters (deposition density,
effective migration time).  The EKI delivers posterior uncertainty bands
without the computational cost of MCMC sampling.

The software is released as the Python package \texttt{soilactivity} on
the \texttt{feature/fredholm-methods} branch, with full test coverage
via synthetic twin experiments.
\end{abstract}

\medskip
\noindent\textbf{Keywords:} Fredholm inversion, in-situ gamma spectrometry,
radionuclide depth profile, total variation, focusing inversion,
ensemble Kalman inversion, advection--dispersion equation, sorption,
depth weighting, geophysical inverse problems

\newpage
% =================================================================
\section{Introduction}\label{sec:intro}

In-situ gamma spectrometry is the primary field technique for rapid
assessment of soil contamination by gamma-emitting radionuclides such as
Cs-137, Cs-134, Eu-152, Co-60, and Am-241.  Unlike destructive soil
core sampling, which provides point measurements at discrete depths, an
in-situ detector placed at height $h$ above the ground records the
integral contribution of all subsurface sources, weighted by the
probability that a photon emitted at depth $z$ reaches the detector
without being absorbed or scattered out of the full-energy peak.  This
measurement geometry leads to a linear inverse problem of a specific
mathematical structure: the relationship between the measured count rates
and the unknown depth-varying volumetric activity $a(z)$ is governed by a
Fredholm integral equation of the first kind
\citep{beck1968deconvolution,zombori1992new}.

The Fredholm formulation was established by \citet{beck1968deconvolution}
and subsequently refined by \citet{zombori1992new}, who demonstrated that
the point-source kernel reduces to an analytic expression involving the
exponential integral $\Eone$ when Compton buildup and angular efficiency
variations are neglected.  The resulting integral equation

\begin{equation}\label{eq:fredholm}
  C_i = \int_0^\infty k_i(z)\, a(z)\, dz, \qquad i = 1,\ldots, m,
\end{equation}

where $C_i$ is the net count rate in the $i$-th gamma line and

\begin{equation}\label{eq:kernel}
  k_i(z) = \frac{y_i\,\varepsilon}{2}\,\Eone\bigl(\mu_{\mathrm{air},i}\,h + \mu_{\mathrm{soil},i}\,z\bigr)
\end{equation}

is the depth kernel, is a classical example of a discrete ill-posed
problem \citep{hansen1998rank}.  The exponential integral $\Eone$ decays
monotonically with depth, so the kernel matrix $K_{ij} = k_i(z_j)\,\Delta z_j$
has rapidly decreasing column norms.  This creates two fundamental
difficulties.  First, small singular values of $K$ amplify measurement
noise, making direct inversion (e.g., via the Moore--Penrose
pseudo-inverse) meaningless without regularisation.  Second, the depth-
dependent sensitivity decay causes regularised solutions to be biased
toward the surface, under-representing activity at depth even when the
data contain sufficient information to resolve it.

These challenges are structurally identical to those encountered in
geophysical inverse problems --- gravity, magnetic, electromagnetic, and
seismic tomography inversions all involve forward operators whose
sensitivity decays with distance from the sensors
\citep{zhdanov2002geophysical,li19963d,li19983d,priede2023linear}.
Over the past two decades, the geophysics community has developed a rich
arsenal of regularisation strategies specifically designed to handle such
problems: total variation (TV) regularisation for piecewise-constant
models \citep{vatankhah2018tv,molodtsov2021convex}, sparsity-promoting
\(\ell_1\) inversion via proximal methods \citep{beck2009fast,boyd2011distributed},
focusing inversion that drives irrelevant model parameters to zero
\citep{portniaguine1999focusing}, depth weighting to counteract sensitivity
decay \citep{li19963d,li19983d,astrakhantsev2021adaptive}, and iterative
methods with implicit regularisation through early stopping
\citep{chen2024adaptive,strohmer2009randomized}.

A second dimension of the problem, largely unexplored in the
radionuclide depth-profiling literature, is the use of physicochemical
transport models as informed priors.  The migration of radionuclides
in soil is governed by the advection--dispersion equation (ADE) with
sorption, decay, and chain ingrowth
\citep{iaea2010radiological,sheppard1990parameter}.  The distribution
coefficient $K_d$ --- which controls how strongly a nuclide is retarded
by sorption to the soil matrix --- depends on soil pH, mineralogy,
organic carbon content, and competing ions
\citep{sheppard1990parameter,roushdy2025sorption}.  For trace
concentrations typical of environmental contamination, the Freundlich
isotherm predicts concentration-dependent $K_d$ that increases at lower
concentrations.  The two-site kinetic sorption model
\citep{vangenuchten1989two,sperry2018two} captures the time-dependent
approach to equilibrium that occurs when only a fraction of sorption
sites are instantaneously accessible.  Incorporating these physics into
the inversion framework yields parametric profiles (pulse, exponential)
with directly interpretable parameters (deposition density, effective
migration time, retardation factor) and provides a natural coupling
between the inversion and the contaminant transport history.

This paper presents the \texttt{soilactivity} software framework, which
bridges these two research streams by implementing a comprehensive suite
of geophysical inversion methods adapted to the Fredholm gamma-spectrometry
kernel, integrated with a reactive transport chemistry module that
provides physically informed priors.  The scientific contributions are:

\begin{enumerate}[label=(\roman*)]
  \item A systematic adaptation of modern geophysical inversion algorithms
        (TV/ADMM, FISTA, IRLS focusing, CGLS, Kaczmarz) to the specific
        structure of the gamma-transport Fredholm equation, including
        depth weighting strategies from gravity/magnetic inversion.
  \item An extended transport-chemistry module with pH-dependent $K_d$,
        Freundlich isotherms, competitive ion exchange, Eh-dependent
        uranium speciation, two-site kinetic sorption, multi-layer soil
columns, and a Crank--Nicolson ADE solver.
  \item A Bayesian uncertainty quantification layer via ensemble Kalman
        inversion (EKI) and Laplace MAP approximation.
  \item An automated pipeline with ensemble multi-method inversion and
        AIC/BIC model selection.
  \item Six regularisation parameter selection criteria (GCV, L-curve,
        discrepancy, quasi-optimality, NCP, SNR) and a comprehensive
        diagnostics suite (resolution matrix, spread function,
        checkerboard test, model covariance).
\end{enumerate}

The remainder of this paper is organised as follows.  Section~2
develops the mathematical formulation.  Section~3 describes the inversion
algorithms.  Section~4 presents the transport-chemistry framework.
Section~5 covers Bayesian methods.  Section~6 details the software
architecture.  Section~7 presents numerical experiments.  Section~8
contains the discussion, and Section~9 the conclusions.


% =================================================================
\section{Mathematical Formulation}\label{sec:math}

\subsection{The Fredholm Equation for Gamma-Transport}

Consider a gamma detector placed at height $h$ above a flat, homogeneous
soil surface.  Let $a(z)$ denote the volumetric activity
$[\text{Bq/m}^3]$ of a given radionuclide at depth $z$ (positive
downwards).  For a gamma line with energy $E$ keV, yield $y$ photons per
disintegration, soil linear attenuation coefficient $\mu_s$ [m$^{-1}$],
and air attenuation coefficient $\mu_a$ [m$^{-1}$], the net count rate
recorded in this line is \citep{beck1968deconvolution}:

\begin{equation}\label{eq:full_kernel}
  C = \frac{y\,\varepsilon}{2}\int_0^{\pi/2}\!
      \tan\theta\; g(\theta)\; B(\tau)\; e^{-\tau}\, d\theta,
\end{equation}

where the optical depth $\tau = (\mu_a h + \mu_s z)/\cos\theta$, $\varepsilon$
is the full-sphere detection efficiency, $g(\theta)$ is the angular
response of the detector, and $B(\tau)$ is the Compton buildup factor
accounting for scattered photons that still deposit full energy in the
detector.

When buildup is neglected ($B=1$) and the angular response is isotropic
($g=1$), the angular integral reduces to the exponential integral:

\begin{equation}\label{eq:analytic}
  k(z) = \frac{y\,\varepsilon}{2}\;\Eone\bigl(\mu_a h + \mu_s z\bigr),
\end{equation}

where $\Eone(x) = \int_1^\infty e^{-xt}/t\, dt$ is the exponential
integral of the first kind.  This analytic form, derived by
\citet{beck1968deconvolution}, is the foundation of all subsequent work.

The kernel depends on the gamma-line energy through $\mu_s(E)$ and
$\mu_a(E)$, which are obtained from NIST databases or empirical
parametrisations.  Table~\ref{tab:lines} lists typical lines and
attenuation coefficients used in the framework.

\begin{table}[htbp]
\centering
\caption{Representative gamma lines and linear attenuation coefficients
(\(\rho_{\mathrm{soil}} = 1.4\)~g/cm$^3$).}
\label{tab:lines}
\begin{tabular}{lcccc}
\toprule
Nuclide & Energy [keV] & Yield & $\mu_s$ [m$^{-1}$] & $\mu_a$ [m$^{-1}$] \\
\midrule
Cs-137 & 661.66 & 0.851 & 10.8 & 0.0092 \\
Eu-152 & 121.78 & 0.284 & 23.8 & 0.0185 \\
Eu-152 & 344.28 & 0.266 & 12.2 & 0.0114 \\
Eu-152 & 1408.01 & 0.209 & 7.4  & 0.0066 \\
Co-60  & 1173.24 & 0.998 & 8.1  & 0.0083 \\
Co-60  & 1332.50 & 0.999 & 7.7  & 0.0078 \\
Am-241 & 59.54  & 0.360 & 44.0 & 0.0300 \\
\bottomrule
\end{tabular}
\end{table}

\subsection{Discretisation and the Ill-Posed Linear System}

Discretising the depth axis into $n$ cells with centres $z_j$ and
widths $\Delta z_j$, and stacking $m$ measurements (multiple lines and/or
detector heights), equation~\eqref{eq:fredholm} becomes:

\begin{equation}\label{eq:linear}
  \bm{d} = \bm{K}\,\bm{a} + \bm{\varepsilon},
\end{equation}

where $\bm{d} \in \R^m$ is the vector of measured count rates,
$\bm{a} \in \R^n$ is the unknown activity profile, $\bm{K} \in \R^{m \times n}$
is the kernel matrix with elements $K_{ij} = k_i(z_j)\,\Delta z_j$, and
$\bm{\varepsilon} \sim \mathcal{N}(\bm{0}, \mathrm{diag}(\sigma^2))$ represents
measurement noise with standard deviations $\sigma_i = \sqrt{C_i}$ (Poisson
statistics).

The singular value decomposition (SVD) of the depth-normalised kernel
$\bm{A} = \bm{K}\,\mathrm{diag}(\bm{s})$ (where $s_j = 1/\|\bm{K}_{\cdot,j}\|$)
reveals the essential ill-posedness of the problem.  For a typical
configuration with $m = 6$ data (three Eu-152 lines at two detector
heights) and $n = 50$ depth cells over $z_{\max} = 0.6$~m, the singular
values span several orders of magnitude.  The ratio of largest to smallest
singular value (the condition number) typically exceeds $10^4$,
confirming severe ill-conditioning.  The Picard plot --- singular values
$\sigma_k$ versus the Fourier coefficients $|\bm{u}_k^\mathrm{T}\bm{d}|$
--- shows the characteristic transition from signal-dominated to
noise-dominated components that is the hallmark of discrete ill-posed
problems \citep{hansen1998rank}.

\subsection{Extended Kernel Models}

Beyond the basic point-source kernel, the framework implements three
extended models that broaden the applicability of the inversion:

\begin{enumerate}[label=\textbf{K\arabic*:}]
  \item \textbf{Lateral extension} ($k_{\mathrm{lat}}$): for wide-area
  surveys (aerial, large ground areas), photons arrive from off-axis
  points within a circular field of view of radius $r_{\mathrm{FOV}}$.
  The lateral kernel integrates the point-source contribution over the
  radial distance $r$ \citep{tyler2008determination}:
  \begin{equation}
    k_{\mathrm{lat}}(z) = \int_0^{r_{\mathrm{FOV}}} k_{\mathrm{point}}(r, z)\; 2\pi r\, dr,
  \end{equation}
  where the path length through soil is $\sqrt{z^2 + r^2}$ and
  through air is $\sqrt{h^2 + r^2}$.

  \item \textbf{Collimated detector} ($k_{\mathrm{coll}}$): for shielded
detectors with a conical field of view of half-angle $\theta_c$,
  the angular integral is truncated at $u_{\max} = 1/\cos\theta_c$
  with a smooth tanh window to avoid numerical Gibbs artifacts:
  \begin{equation}
    k_{\mathrm{coll}}(z) = \frac{y\,\varepsilon}{2}
      \int_1^{u_{\max}} \frac{B(\tau)\,e^{-\tau}\,w(u)}{u}\, du,
  \end{equation}
  where $w(u) = \tfrac{1}{2}(1 - \tanh(20(u - u_{\max})))$.

  \item \textbf{Multi-layer soil} ($k_{\mathrm{ml}}$): for stratified
  soils where each layer has a different density and composition, the
  attenuation coefficient $\mu_s$ is a step function of depth, and the
  kernel is evaluated with the layer-appropriate $\mu_s(z_j)$ at each
  grid point.
\end{enumerate}

\subsection{Buildup Factor}

The Taylor buildup factor $B(\tau) = 1 + a\,\tau\,e^{-\tau/b}$ provides a
simple parametric approximation to Compton scattering contributions.
With default parameters $a = 1.0$, $b = 8.0$, it captures the
characteristic rise-and-plateau behaviour of the buildup factor as a
function of optical depth.  More sophisticated parametrisations (e.g.,
the Geometric Progression formula) can be supplied as callable arguments
to the kernel construction function.


% =================================================================
\section{Inversion Algorithms}\label{sec:solvers}

All solvers in the framework operate on the weighted system
$\bm{A} = \mathrm{diag}(\bm{w})\,\bm{K}$, $\bm{b} = \bm{w}\,\bm{d}$,
where $w_i = 1/\sigma_i$, and enforce non-negativity ($a(z) \geq 0$) by
default, reflecting the physical constraint that volumetric activity
cannot be negative.

\subsection{Classical Solvers}

\subsubsection{Tikhonov Regularisation}

The Tikhonov solution solves the constrained optimisation problem:

\begin{equation}\label{eq:tikhonov}
  \bm{a}_\alpha = \argmin_{\bm{a} \geq 0}\;
    \|\bm{A}\bm{a} - \bm{b}\|^2 + \alpha\,\|\bm{L}(\bm{a} - \bm{a}_0)\|^2,
\end{equation}

where $\bm{L}$ is the regularisation operator (identity for minimum-norm,
second-order finite-difference matrix for smoothness), $\bm{a}_0$ is a
reference model (zero by default), and $\alpha > 0$ is the regularisation
parameter.  The problem is solved as a non-negative least squares (NNLS)
problem via \texttt{scipy.optimize.lsq\_linear} after augmenting the system:

\begin{equation}
  \begin{pmatrix} \bm{A} \\ \sqrt{\alpha}\,\bm{L} \end{pmatrix}
  \bm{a} = \begin{pmatrix} \bm{b} \\ \sqrt{\alpha}\,\bm{L}\,\bm{a}_0 \end{pmatrix},
  \quad \bm{a} \geq 0.
\end{equation}

The smoothness prior ($\bm{L} = \bm{D}_2$, the second-order finite-difference
matrix) produces smoothly varying profiles suitable for diffusively
migrated nuclides.  The minimum-norm prior ($\bm{L} = \bm{I}$) is used as
a baseline for comparison with more sophisticated methods.

\subsubsection{Truncated SVD (TSVD)}

TSVD applies the depth-normalised SVD, $\bm{A}\,\bm{S} = \bm{U}\,\bm{\Sigma}\,\bm{V}^\mathrm{T}$,
and retains only the $k$ largest singular values:

\begin{equation}
  \bm{a}_{\mathrm{TSVD}} = \bm{S}\,\sum_{i=1}^{k}
  \frac{\bm{u}_i^\mathrm{T}\bm{b}}{\sigma_i}\,\bm{v}_i,
\end{equation}

where the truncation level $k$ is determined by the relative cutoff
$\sigma_k > \sigma_1 \times 10^{-3}$.  The depth normalisation
$\bm{S} = \mathrm{diag}(1/\|\bm{K}_{\cdot,j}\|)$ ensures that each depth
layer contributes equally to the SVD, counteracting the natural dominance
of near-surface layers \citep{li19963d}.

\subsubsection{Landweber Iteration}

The Landweber method is a steepest-descent iteration on the normal
equations with relaxation parameter $\omega$:

\begin{equation}
  \bm{a}^{(k+1)} = \max\!\bigl(\bm{a}^{(k)} + \omega\,\bm{A}^\mathrm{T}(\bm{b} - \bm{A}\bm{a}^{(k)}),\; 0\bigr),
\end{equation}

where $\omega = 0.9 \times 2/\sigma_1^2$ ensures convergence.  The iteration
count $k$ acts as the implicit regularisation parameter through the
phenomenon of semi-convergence: the solution first improves (approaching
the true profile) and then deteriorates (noise amplification) as
$k$ increases.  Stopping at the discrepancy principle
$\|\bm{A}\bm{a}^{(k)} - \bm{b}\|^2 \leq m$ prevents overfitting.

\subsubsection{Cimmino (Simultaneous Kaczmarz)}

The Cimmino method projects onto all equation hyperplanes simultaneously
at each iteration:

\begin{equation}
  \bm{a}^{(k+1)} = \bm{a}^{(k)} + \frac{1}{m}
    \sum_{i=1}^{m} \frac{r_i^{(k)}}{\|\bm{A}_{i\cdot}\|^2}\,\bm{A}_{i\cdot},
\end{equation}

where $r_i^{(k)} = b_i - \bm{A}_{i\cdot}\,\bm{a}^{(k)}$ is the $i$-th residual.
Cimmino is particularly convenient for the block-structured systems that
arise from multi-line, multi-height measurements, as it naturally
handles the row partitioning \citep{hansen1998rank}.

\subsection{Geophysical Solvers}

\subsubsection{Conjugate Gradient Least Squares (CGLS)}

CGLS applies the conjugate gradient algorithm to the normal equations
$\bm{A}^\mathrm{T}\bm{A}\,\bm{a} = \bm{A}^\mathrm{T}\bm{b}$, providing $\mathcal{O}(1/k)$
convergence to the minimum-norm solution.  Following
\citet{chen2024adaptive}, the iteration count $k$ serves as the
regularisation parameter, with the L-curve corner used for automatic
stopping.  The algorithm maintains orthogonal search directions and
requires only matrix-vector products, making it efficient for large
discretisations.  The non-negativity constraint is enforced by projecting
onto $\bm{a} \geq 0$ after each update.

\subsubsection{Randomised Kaczmarz}

The Algebraic Reconstruction Technique (ART / Kaczmarz) cycles through
the equations sequentially, projecting onto each hyperplane:

\begin{equation}
  \bm{a}^{(k+1)} = \bm{a}^{(k)} + \frac{b_i - \bm{A}_{i\cdot}\,\bm{a}^{(k)}}{\|\bm{A}_{i\cdot}\|^2}\,\bm{A}_{i\cdot},
\end{equation}

where the row index $i$ is selected randomly with replacement (without
replacement within each cycle).  \citet{strohmer2009randomized} proved
that randomised ordering achieves exponential convergence for
row-norm-normalised systems, in contrast to the slow convergence of
sequential cycling for ill-conditioned problems.

\subsubsection{FISTA: $\ell_1$-Sparse Inversion}

FISTA \citep{beck2009fast} solves the $\ell_1$-regularised problem:

\begin{equation}\label{eq:fista}
  \bm{a}^* = \argmin_{\bm{a} \geq 0}\; \tfrac{1}{2}\|\bm{A}\bm{a} - \bm{b}\|^2 + \alpha\,\|\bm{a}\|_1.
\end{equation}

The $\ell_1$ penalty promotes sparsity, concentrating activity in a
few thin layers --- ideal for compact contamination plumes.  FISTA uses
Nesterov acceleration, achieving $\mathcal{O}(1/k^2)$ convergence versus
$\mathcal{O}(1/k)$ for standard ISTA.  Each iteration consists of:

\begin{enumerate}
  \item Gradient step: $\tilde{\bm{a}} = \bm{y} - \tfrac{1}{L}\,\bm{A}^\mathrm{T}(\bm{A}\bm{y} - \bm{b})$,
        where $L = \sigma_1^2$ is the Lipschitz constant.
  \item Proximal step (non-negative soft thresholding):
        $a_j = \max(\tilde{a}_j - \alpha/L,\; 0)$.
  \item Momentum update: $\bm{y} \leftarrow \bm{a} + \frac{t-1}{t_{\mathrm{new}}}(\bm{a} - \bm{a}_{\mathrm{old}})$,
        $t_{\mathrm{new}} = \frac{1}{2}(1 + \sqrt{1 + 4t^2})$.
\end{enumerate}

\subsubsection{Total Variation via ADMM}

Total variation regularisation preserves sharp layer boundaries by
penalising the magnitude of gradients rather than the gradients
themselves:

\begin{equation}\label{eq:tv}
  \bm{a}^* = \argmin_{\bm{a} \geq 0}\; \tfrac{1}{2}\|\bm{A}\bm{a} - \bm{b}\|^2 + \alpha\,\mathrm{TV}(\bm{a}),
\end{equation}

where $\mathrm{TV}(\bm{a}) = \sum_j |(\bm{D}\bm{a})_j|$ for the anisotropic
case ($\bm{D}$ is the first-order forward-difference operator).  The
framework uses the ADMM (Alternating Direction Method of Multipliers)
\citep{boyd2011distributed} formulation with the split variable
$\bm{z} = \bm{D}\bm{a}$:

\begin{align}
  \text{\textbf{m-update:}} &\quad
    (\bm{A}^\mathrm{T}\bm{A} + \rho\,\bm{D}^\mathrm{T}\bm{D})\,\bm{a} = \bm{A}^\mathrm{T}\bm{b} + \rho\,\bm{D}^\mathrm{T}(\bm{z} - \bm{u}), \label{eq:admm_m}\\
  \text{\textbf{z-update:}} &\quad
    \bm{z} = \mathcal{S}_{\alpha/\rho}(\bm{D}\bm{a} + \bm{u}), \label{eq:admm_z}\\
  \text{\textbf{u-update:}} &\quad
    \bm{u} \leftarrow \bm{u} + \bm{D}\bm{a} - \bm{z}, \label{eq:admm_u}
\end{align}

where $\mathcal{S}_\tau$ is the soft-thresholding operator and $\rho$ is
the ADMM penalty parameter.  The system matrix in
equation~\eqref{eq:admm_m} is constant throughout the iterations and is
prefactored via Cholesky decomposition, making each m-update an
$\mathcal{O}(n)$ back-substitution.  This formulation follows
\citet{vatankhah2018tv}, adapted from 3D gravity to 1D depth profiling.
TV regularisation is particularly well-suited for contamination
scenarios with abrupt layer boundaries (e.g., a buried waste container
or a distinct plume interface).

\subsubsection{Focusing Inversion (IRLS)}

The focusing (compact) inversion of \citet{portniaguine1999focusing}
minimises a stabilising functional that drives most model parameters to
zero, producing compact, sharp-boundary solutions:

\begin{equation}
  \bm{a}^* = \argmin_{\bm{a} \geq 0}\; \|\bm{A}\bm{a} - \bm{b}\|^2 + \alpha\, S(\bm{a}),
\end{equation}

where the stabiliser $S$ takes one of two forms.  The Minimum Support (MS)
stabiliser penalises non-zero values:

\begin{equation}
  S_{\mathrm{MS}}(\bm{a}) = \sum_j \frac{a_j^2}{a_j^2 + e^2},
\end{equation}

producing sparse solutions where only a few depth layers have
non-negligible activity.  The Minimum Gradient Support (MGS) stabiliser
penalises non-zero gradients:

\begin{equation}
  S_{\mathrm{MGS}}(\bm{a}) = \sum_j \frac{(\bm{D}\bm{a})_j^2}{(\bm{D}\bm{a})_j^2 + e^2},
\end{equation}

producing blocky, layered models with sharp interfaces and flat
interiors --- precisely the structure expected for stratified
contamination.  The parameter $e$ controls the sharpness of boundaries
(smaller $e$ = sharper but potentially unstable) and is progressively
decreased during the IRLS iterations following the prescription of
\citet{zhdanov2002geophysical}.  At each IRLS iteration, the
non-quadratic stabiliser is locally approximated by a quadratic form
with reweighted diagonal matrices, reducing the problem to a sequence
of Tikhonov-like systems:

\begin{equation}
  \bigl(\bm{A}^\mathrm{T}\bm{A} + \alpha\,\bm{W}^{-1}\bigr)\,\bm{a} = \bm{A}^\mathrm{T}\bm{b},
\end{equation}

where $\bm{W}$ is the diagonal matrix of weights constructed from the
current solution.

\subsection{Regularisation Parameter Selection}\label{sec:criteria}

Six parameter selection criteria are implemented, providing redundant
and complementary approaches to the fundamental problem of balancing
data fidelity against solution stability.

\subsubsection{Generalised Cross-Validation (GCV)}

GCV \citep{golub1979generalized} minimises:

\begin{equation}
  \mathrm{GCV}(\alpha) = \frac{m\,\|\bm{A}\bm{a}_\alpha - \bm{b}\|^2}{\bigl[\tr(\bm{I} - \bm{H}_\alpha)\bigr]^2},
\end{equation}

where $\bm{H}_\alpha = \bm{A}(\bm{A}^\mathrm{T}\bm{A} + \alpha\,\bm{L}^\mathrm{T}\bm{L})^{-1}\bm{A}^\mathrm{T}$
is the hat (influence) matrix.  The trace is estimated via Hutchinson's
stochastic trace estimator with random $\pm 1$ probe vectors, avoiding
the $\mathcal{O}(n^3)$ cost of explicit computation.  A weighted GCV
variant handles heteroscedastic (Poisson) noise by incorporating
$\sigma_i$ into the weighting.

\subsubsection{L-Curve Corner}

The L-curve \citep{hansen1992analysis} is the parametric plot of
$(\log\|\bm{L}\bm{a}_\alpha\|,\, \log\|\bm{A}\bm{a}_\alpha - \bm{b}\|)$
as $\alpha$ varies.  The optimal $\alpha$ is at the corner of maximum
curvature, computed via the Menger curvature formula for three
consecutive points in log-log space.  An iterative variant
$\texttt{lcurve\_corner\_iter}$ applies the same curvature criterion to
the residual/solution-norm trajectories of iterative methods (CGLS,
Landweber).

\subsubsection{Discrepancy Principle}

The Morozov discrepancy principle \citep{morozov1966choice} selects the
largest $\alpha$ such that the weighted residual equals the expected
value under the noise model:

\begin{equation}
  \|\bm{A}\bm{a}_\alpha - \bm{b}\|^2 = m.
\end{equation}

This is implemented via Brent's root-finding method on a log-spaced
grid of $\alpha$ values.

\subsubsection{Quasi-Optimality Criterion}

The quasi-optimality criterion \citep{hochstenbach2015filtering}
selects $\alpha$ minimising the noise-dominated SVD components:

\begin{equation}
  \alpha^* = \argmin_\alpha \sqrt{\sum_{k:\,\sigma_k^2 < \alpha}
    \left(\frac{\bm{u}_k^\mathrm{T}\bm{b}}{\sigma_k}\right)^2}.
\end{equation}

This criterion does not require a noise-level estimate, making it useful
when $\sigma_i$ is uncertain.

\subsubsection{Normalised Cumulative Periodogram (NCP)}

The NCP criterion tests whether the residuals are consistent with white
noise (as expected for the correct regularisation level) using the
Kolmogorov--Smirnov statistic on the cumulative periodogram of the
residuals.  The optimal $\alpha$ minimises the KS distance, indicating
the most white-noise-like residual.

\subsubsection{Signal-to-Noise Ratio (SNR) Criterion}

The SNR criterion maximises the estimated signal-to-noise ratio of the
solution:

\begin{equation}
  \mathrm{SNR}(\alpha) = \frac{\|\bm{a}_\alpha\|^2}{\sigma^2\,\tr(\bm{C}_{\mathrm{post}})} = \frac{\sum_k (\sigma_k\,\beta_k)^2/(\sigma_k^2 + \alpha)^2}{\sum_k 1/(\sigma_k^2 + \alpha)},
\end{equation}

where $\beta_k = \bm{u}_k^\mathrm{T}\bm{b}$ and $\sigma^2$ is estimated from the
residuals.


% =================================================================
\section{Transport Chemistry}\label{sec:transport}

\subsection{Advection--Dispersion Equation with Sorption}

The one-dimensional ADE governing the vertical migration of a
radionuclide in a saturated soil column is:

\begin{equation}\label{eq:ade}
  R\,\frac{\partial C}{\partial t} = \frac{\partial}{\partial z}\!
    \left(D\,\frac{\partial C}{\partial z}\right) - v\,\frac{\partial C}{\partial z}
    - \lambda\,R\,C,
\end{equation}

where $C(z,t)$ is the dissolved-phase concentration [Bq/L],
$D$ is the hydrodynamic dispersion coefficient [m$^2$/s],
$v$ is the pore-water velocity [m/s] (positive downward),
$\lambda$ is the radioactive decay constant [s$^{-1}$], and
$R$ is the retardation factor:

\begin{equation}\label{eq:retardation}
  R = 1 + \frac{\rho_b\,K_d}{\theta},
\end{equation}

with $\rho_b$ the bulk density [g/cm$^3$], $K_d$ the distribution
coefficient [mL/g], and $\theta$ the volumetric water content.
For Cs-137 in typical mineral soils, $R$ ranges from 100 to 5000,
meaning that the nuclide migrates 100--5000 times slower than the pore
water.

\subsection{Sorption Models}

\subsubsection{Linear $K_d$ Database}

The framework includes a comprehensive $K_d$ database covering 12
radionuclides (Cs-137, Sr-90, Co-60, U-238, Pu-239, Am-241, Ra-226,
Pb-210, Eu-152, I-131, Tc-99, Se-79) with low/mid/high scenarios
based on IAEA TRS-472 \citep{iaea2010radiological} and Sheppard \\
\& Thibault \citep{sheppard1990parameter}.  The geometric mean of the
low--high range provides the ``mid'' estimate used for default
parameterisation.

\subsubsection{pH-Dependent $K_d$}

For the most critical nuclides, the framework implements a
dependent $K_d$ model:

\begin{equation}
  \log_{10} K_d = a_{\mathrm{pH}}\,\mathrm{pH} + b_{\mathrm{clay}}\,f_{\mathrm{clay}} + c_{\mathrm{OC}}\,f_{\mathrm{OC}} + d,
\end{equation}

where the coefficients are fitted to the IAEA compilation data.  This
allows site-specific $K_d$ estimation from routinely measured soil
parameters (pH, clay fraction, organic carbon fraction).

\subsubsection{Freundlich Isotherm}

At trace concentrations typical of environmental contamination, the
linear $K_d$ model may underestimate sorption because the Freundlich
isotherm predicts:

\begin{equation}
  K_d(C) = K_F\,C^{n_F - 1},
\end{equation}

where $K_F$ [mL/g $\cdot$ (L/g)$^{n_F-1}$] and $n_F < 1$ (favourable
isotherm).  As $C \to 0$, $K_d \to \infty$, reflecting the extremely
strong sorption of radionuclides at trace levels.  The framework
includes Freundlich parameters for Cs-137, Sr-90, Pu-239, and U-238.

\subsubsection{Competitive Ion Exchange}

Cs$^+$ competes with K$^+$ and NH$_4^+$ for frayed-edge sites (FES) on
illite and vermiculite clay minerals:

\begin{equation}
  K_{d,\mathrm{Cs}}^{\mathrm{eff}} = \frac{K_{d,\mathrm{Cs}}^{\mathrm{ref}}}{1 + \alpha_{\mathrm{K}}\,[\mathrm{K}^+] + \alpha_{\mathrm{NH}_4}\,[\mathrm{NH}_4^+]},
\end{equation}

where $\alpha_{\mathrm{K}} = 0.1$ and $\alpha_{\mathrm{NH}_4} = 0.4$ are
selectivity coefficients.  Similarly, Sr$^{2+}$ competes with
Ca$^{2+}$: $K_{d,\mathrm{Sr}}^{\mathrm{eff}} = K_{d,\mathrm{Sr}}^{\mathrm{ref}} / (1 + \alpha_{\mathrm{Ca}}\,[\mathrm{Ca}^{2+}])$.

\subsubsection{Eh-Dependent Uranium Speciation}

Uranium migration is controlled by the redox potential through the
U(VI)/U(IV) transition.  Under oxidising conditions, U(VI) forms
mobile uranyl-carbonate complexes; under reducing conditions, U(IV)
precipitates as immobile uraninite.  The framework implements:

\begin{equation}
  \log_{10} K_d = \begin{cases}
    2.0 + 0.3\,\mathrm{pH}  & \text{if } E_h < E_h^{\mathrm{crit}}, \\
    -0.5 + 0.15\,\mathrm{pH} - 0.1\,\log_{10}[\mathrm{HCO}_3^-] & \text{otherwise,}
  \end{cases}
\end{equation}

where $E_h^{\mathrm{crit}} = 470 - 60\,\mathrm{pH}$ mV.

\subsection{Two-Site Kinetic Sorption}

The two-site model \citep{vangenuchten1989two} partitions sorption
sites into an instantaneous (equilibrium) fraction $f$ and a kinetic
fraction $1 - f$ governed by first-order rate exchange.  The effective
retardation factor becomes time-dependent:

\begin{equation}
  R_{\mathrm{eff}}(t) = 1 + \frac{\rho_b\,K_d}{\theta}\n    \bigl[f + (1-f)(1 - e^{-\omega t})\bigr],
\end{equation}

where $\omega$ is the kinetic rate constant [s$^{-1}$].  At short times
($t \ll 1/\omega$), only equilibrium sites contribute; at long times,
the full $K_d$ is recovered.  The framework implements a coupled
ODE system for the dissolved and kinetically-sorbed phases, solved via
operator splitting (Crank--Nicolson for transport, analytical integration
for kinetic exchange).

\subsection{Multi-Layer Soil Columns}

For stratified soils, the ADE is solved independently in each layer with
its own $D$, $v$, $R$, and $\lambda$, with continuity of flux and
concentration enforced at layer interfaces.  This allows modelling of
typical field scenarios such as a low-permeability clay layer beneath a
more permeable topsoil, or a buried contaminated layer beneath clean
backfill.

\subsection{Analytical Profile Functions}

Two analytical solutions of the ADE serve as parametric profile models
for the inversion:

\textbf{Pulse profile} (instantaneous surface deposition $A_0$ at $t = 0$):
\begin{equation}
  a(z,t) = \frac{A_0\,e^{-\lambda t}}{\sqrt{4\pi D_s t}}\n    \Bigl[e^{-\frac{(z - v_s t)^2}{4 D_s t}} + e^{-\frac{(z + v_s t)^2}{4 D_s t}}\Bigr],
\end{equation}

where $D_s = D/R$ and $v_s = v/R$.  The image source at $-z$ enforces
the no-flux boundary condition at the surface.

\textbf{Exponential profile} (diffusion-limited relaxation):
\begin{equation}
  a(z) = \frac{A}{\lambda_{\mathrm{relax}}}\,e^{-z/\lambda_{\mathrm{relax}}},
\end{equation}

normalised so that $\int_0^\infty a(z)\,dz = A$ [Bq/m$^2$].

\subsection{Bateman Decay Chains}

The matrix exponential formulation of the Bateman equations handles
arbitrary decay chains:

\begin{equation}
  \bm{A}(t) = \exp(\bm{M}\,t)\,\bm{A}(0),
\end{equation}

where $M_{ii} = -\lambda_i$ and $M_{i,i-1} = \lambda_{i-1}\,b_{i-1}$
(branching ratio).  This allows modelling of ingrowth from parent
daughters (e.g., Cs-137 $\to$ Ba-137m) within the inversion framework.

\subsection{Crank--Nicolson Numerical Solver}

For problems where analytical solutions are unavailable (time-dependent
boundary conditions, spatially varying parameters), the framework
provides a Crank--Nicolson finite-difference solver for the ADE.  The
scheme is unconditionally stable and second-order accurate in both
space and time, using the Thomas algorithm for the resulting tridiagonal
system at each time step.  Harmonic-mean interface diffusivities ensure
conservation of flux at layer boundaries.


% =================================================================
\section{Bayesian Methods}\label{sec:bayesian}

\subsection{Ensemble Kalman Inversion (EKI)}

EKI \citep{iglesias2013ensemble,laby2025regularised,tso2024ensemble}
approximates the Bayesian posterior $p(\bm{a} \mid \bm{d})$ using an
ensemble of model realisations updated via Kalman filter equations.
Given an ensemble $\{\bm{a}_j\}_{j=1}^{N_e}$ initialised from the prior
$\mathcal{N}(\bm{a}_0, \mathrm{diag}(\sigma_p^2))$, each iteration
consists of:

\begin{enumerate}
  \item \textbf{Forward step:} compute synthetic data for each member,
        $\bm{d}_j = \bm{K}\,\bm{a}_j$.
  \item \textbf{Statistics:} compute ensemble means and cross-covariances
        $\bm{C}_{ma}$, $\bm{C}_d$.
  \item \textbf{Kalman update:}
        $\bm{a}_j \leftarrow \bm{a}_j + \bm{C}_{ma}\,\bm{C}_d^{-1}(\bm{d}_{\mathrm{obs}} + \bm{\eta}_j - \bm{d}_j)$,
        where $\bm{\eta}_j \sim \mathcal{N}(\bm{0}, \bm{C}_{dd})$ provides
        inflation to prevent ensemble collapse.
  \item \textbf{Non-negativity projection:} $\bm{a}_j \leftarrow \max(\bm{a}_j, 0)$.
\end{enumerate}

The regularised variant of \citet{laby2025regularised} adds
$\alpha\,\bm{I}$ to $\bm{C}_d$ and $\bm{C}_{ma}$ to further prevent collapse.
For linear forward models (as in our case), EKI converges in
approximately 10--15 iterations, providing posterior mean, standard
deviation, and full ensemble for uncertainty quantification --- at a
cost of seconds rather than the hours required for MCMC.

\subsection{Laplace MAP Approximation}

For a Gaussian prior $p(\bm{a}) \sim \mathcal{N}(\bm{a}_0,
(\alpha\,\bm{L}^\mathrm{T}\bm{L})^{-1})$ and Gaussian likelihood
$p(\bm{d} \mid \bm{a}) \sim \mathcal{N}(\bm{K}\bm{a},\,\mathrm{diag}(\sigma^2))$,
the MAP estimate coincides with the Tikhonov solution:

\begin{equation}
  \bm{a}_{\mathrm{MAP}} = (\bm{A}^\mathrm{T}\bm{A} + \alpha\,\bm{L}^\mathrm{T}\bm{L})^{-1}
  \bm{A}^\mathrm{T}\bm{b},
\end{equation}

and the posterior covariance is:

\begin{equation}
  \bm{C}_{\mathrm{post}} = (\bm{A}^\mathrm{T}\bm{A} + \alpha\,\bm{L}^\mathrm{T}\bm{L})^{-1}.
\end{equation}

The diagonal of $\bm{C}_{\mathrm{post}}$ provides 1-$\sigma$ uncertainty
estimates for each depth layer.  This is computationally cheap (one
matrix inversion) but assumes Gaussian posterior, which may
underestimate uncertainty for non-negative, sparse, or multi-modal
posteriors.  A Gaussian Process (RBF) prior covariance is also available
as an alternative to the Tikhonov smoothness prior:

\begin{equation}
  C_{ij} = \sigma_f^2\,\exp\!\left(-\frac{(z_i - z_j)^2}{2\ell^2}\right) + \sigma_n^2\,\delta_{ij},
\end{equation}

encoding smoothness with correlation length $\ell$.


% =================================================================
\section{Depth Weighting and Joint Inversion}\label{sec:geophysics}

\subsection{Depth Weighting Strategies}

The exponential decay of the gamma kernel with depth creates a strong
bias in regularised inversions toward near-surface activity.  Depth
weighting counteracts this by rebalancing the sensitivity of the
inversion across the depth range.  The framework implements four
strategies adapted from the geophysical inversion literature:

\begin{enumerate}[label=\textbf{W\arabic*:}]
  \item \textbf{Li--Oldenburg} (column-norm normalisation) \citep{li19963d}:
  $w_j = 1/\|\bm{K}_{\cdot,j}\|$.  Equalises the column norms of the
  kernel matrix, ensuring that each depth layer contributes equally to
  the normal equations.  This is the standard approach in gravity and
  magnetic inversion.

  \item \textbf{Power-law} \citep{li19983d}:
  $w(z) = (z/z_{\mathrm{ref}} + 1)^{-\beta/2}$.
  Originally developed for 3D gravity inversion where the kernel
  decays as $1/z^3$, adapted here with $\beta \in [1, 3]$ for the
  exponential gamma kernel.

  \item \textbf{Logarithmic}: $w(z) = 1/\log(1 + z/z_0)$.
  Provides milder weighting than power-law, useful when the kernel
  already includes some depth normalisation.

  \item \textbf{Adaptive} (data-driven): combines kernel column norms,
  power-law decay, and a smoothness-based correction that down-weights
  depths where the kernel sensitivity oscillates (indicating poor
  resolution):
  $w_{\mathrm{adapt}} = w_{\mathrm{kern}} \cdot w_{\mathrm{power}} \cdot (1 + \alpha / (|\nabla w_{\mathrm{kern}}| + \epsilon))$.
\end{enumerate}

A depth-weighted Tikhonov solve applies the weighting through the
regularisation operator: $\bm{L}_w = \mathrm{diag}(\sqrt{\bm{w}})\,\bm{L}$.

\subsection{Joint Multi-Nuclide Inversion}

When multiple nuclides are measured simultaneously (e.g., a fallout
mixture of Cs-137, Cs-134, and Eu-152), the framework constructs a
block-diagonal joint kernel:

\begin{equation}
  \bm{K}_{\mathrm{joint}} = \begin{pmatrix}
    \bm{K}_{\mathrm{Cs}} & \bm{0} & \bm{0} \\
    \bm{0} & \bm{K}_{\mathrm{Eu}} & \bm{0} \\
    \bm{0} & \bm{0} & \bm{K}_{\mathrm{Am}}
  \end{pmatrix},
\end{equation}

and solves the stacked system with shared depth grid.  Coupling
constraints (``same shape'', ``proportional'') exploit the fact that
co-deposited nuclides from a common source should share the same depth
profile shape, differing only in amplitude.  This reduces the effective
number of unknowns and improves resolution.

\subsection{Ensemble Multi-Method Inversion}

The pipeline's ensemble mode runs all available solvers
(Tikhonov/GCV, TV/ADMM, FISTA, focusing/MGS, CGLS) on the same data
and selects the best result via information criteria:

\begin{align}
  \mathrm{AIC} &= \chi^2 + 2\,n_{\mathrm{eff}}, \\
  \mathrm{BIC} &= \chi^2 + \ln(m)\,n_{\mathrm{eff}},
\end{align}

where $n_{\mathrm{eff}} = \min(\sum_j \mathbb{1}[a_j > 0.001\,\max(\bm{a})],\, m)$
is the effective number of parameters (number of active depth cells).
BIC penalises model complexity more strongly than AIC, favouring
sparser solutions when data are limited.


% =================================================================
\section{Resolution Diagnostics}\label{sec:diagnostics}

The framework provides a suite of resolution diagnostics adapted from
seismic tomography \citep{backus1968numerical,menke2018geophysical}.

\subsection{Resolution Matrix}

The model resolution matrix $\bm{R} = (\bm{A}^\mathrm{T}\bm{A} + \alpha\,\bm{L}^\mathrm{T}\bm{L})^{-1}
\bm{A}^\mathrm{T}\bm{A}$ shows how a delta-function layer at depth $z_j$ is
recovered across all depths.  Diagonal dominance indicates good vertical
resolution; off-diagonal spread indicates smearing.  The column $\bm{R}_{\cdot,j}$
is the point-spread function (PSF) for depth $z_j$.

\subsection{Depth of Investigation}

The depth of investigation $z^*$ is defined as the depth at which the
cumulative absolute PSF reaches a specified fraction $f$ (default 0.5):

\begin{equation}
  z^*_j = \inf\{z : \sum_{i: z_i \leq z} |R_{ij}| \geq f\}.
\end{equation}

This provides a depth-dependent measure of how deep the measurement
configuration can reliably resolve activity.

\subsection{Spread Function}

The Backus--Gilbert spread function quantifies the spatial spreading
of the PSF:

\begin{equation}
  s_j = 12 \sum_i \bar{R}_{ij}^2\,(z_i - z_j)^2\,\Delta z,
\end{equation}

where $\bar{R}_{ij} = R_{ij}/\sum_k R_{ik}$ is the row-normalised resolution
matrix.  Smaller spread indicates better localisation.

\subsection{Model Covariance}

The linearised posterior model covariance
$\bm{C}_m = (\bm{A}^\mathrm{T}\bm{A} + \alpha\,\bm{L}^\mathrm{T}\bm{L})^{-1}$
provides depth-dependent uncertainty estimates.  The diagonal elements
give 1-$\sigma$ uncertainties for each layer; off-diagonal elements
reveal correlated errors between layers.

\subsection{Additional Diagnostics}

The framework also implements:
(1)~a checkerboard resolution test that measures how well alternating
zero/unity blocks in the depth profile are recovered;
(2)~a data resolution matrix (hat matrix) showing which data points
contribute most to each model parameter;
(3)~per-channel sensitivity kernels showing the depth sensitivity of
each measurement;
(4)~information content metrics (trace of resolution matrix, effective
rank, condition number).


% =================================================================
\section{Software Architecture}\label{sec:software}

The \texttt{soilactivity} package is organised into seven modules within
the \texttt{depth\_inversion} subpackage, following a clear separation
of concerns:

\begin{enumerate}
  \item \textbf{kernels.py}: Forward Fredholm kernel construction
  (analytic, numerical, lateral, collimated, multi-layer) and the
  \texttt{GammaLine}/\texttt{Detector} data classes.
  \item \textbf{transport.py}: Reactive transport chemistry ($K_d$
  database, pH-dependent and Freundlich sorption, competitive ion
  exchange, Eh-dependent U, two-site kinetic sorption, multi-layer
  ADE, Crank--Nicolson solver, Bateman chains).
  \item \textbf{solvers.py}: Ten inversion algorithms (Tikhonov, TSVD,
  Landweber, Cimmino, CGLS, Kaczmarz, FISTA, TV/ADMM, focusing MGS/MS,
  CGLS+L-curve, cross-validation, depth-weighted Tikhonov).
  \item \textbf{criteria.py}: Six parameter selection criteria (GCV,
  L-curve, discrepancy, quasi-optimality, NCP, SNR, weighted GCV).
  \item \textbf{diagnostics.py}: Resolution matrix, DOI, spread function,
  model covariance, checkerboard test, sensitivity kernels, information
  content.
  \item \textbf{bayesian.py}: EKI, Laplace MAP, GP prior.
  \item \textbf{geophysics.py}: Depth weighting (Li--Oldenburg, power-law,
  logarithmic, adaptive), joint multi-nuclide inversion, two-site ADE,
  regularisation operators.
  \item \textbf{pipeline.py}: \texttt{DepthInverter} class providing
  end-to-end inversion with automatic parameter selection, ensemble
  multi-method with AIC/BIC, and parametric transport-informed
  inversion.
\end{enumerate}

The pipeline's \texttt{DepthInverter} class provides a high-level
interface that constructs the kernel matrix from gamma lines and
detector heights, then offers methods for each inversion strategy.
The \texttt{fit()} method performs non-parametric Tikhonov with automatic
$\alpha$ selection; \texttt{fit\_tv()}, \texttt{fit\_sparse()},
\texttt{fit\_focusing()} provide the geophysical alternatives;
\texttt{fit\_eki()} provides Bayesian uncertainty quantification;
\texttt{fit\_ensemble()} runs all methods and selects the best via AIC;
and \texttt{fit\_parametric()} inverts for physically interpretable transport
parameters.  All methods return a \texttt{DepthInversion} dataclass with
the reconstructed profile, uncertainty, areal activity, median depth,
and chi-squared statistic.

The test suite consists of 25+ twin experiments covering kernel
validation (numeric--analytic agreement to 0.2\%), transport chemistry
($K_d$ lookups, retardation, pH dependence, Freundlich, multi-layer,
competitive sorption), all solvers (non-negativity, convergence), all
criteria (finite $\alpha$), diagnostics (resolution, covariance, spread,
information content), Bayesian methods (EKI, Laplace), and the pipeline
ensemble.


% =================================================================
\section{Numerical Experiments}\label{sec:experiments}

\subsection{Twin Experiment Setup}

The twin experiments follow a controlled protocol: (1)~generate a
synthetic activity profile $a_{\mathrm{true}}(z)$; (2)~compute
deterministic count rates $\bm{c} = \bm{K}\,\bm{a}_{\mathrm{true}}$;
(3)~draw Poisson-realised noisy counts $d_i \sim \mathrm{Poisson}(c_i \times\mathit{scale}) / \mathit{scale}$;
(4)~invert to obtain $\hat{\bm{a}}$; (5)~compare with $a_{\mathrm{true}}$.

Two scenarios are used:

\textbf{Scenario A (Eu-152, multi-line):} Three gamma lines at
121.78, 344.28, and 1408.01 keV, measured at two detector heights
(0.5 and 2.0~m), yielding $m = 6$ data for $n = 20$ depth cells over
$z_{\max} = 0.6$~m.  The true profile is a pulse profile with
$A_0 = 10^5$~Bq/m$^2$, $t = 25$~years, $D = 10^{-10}$~m$^2$/s,
$R = 1401$ (Cs-like retardation).

\textbf{Scenario B (Cs-137, single-line):} One line at 661.66 keV,
measured at 1.0~m height, yielding $m = 1$ data for $n = 50$ depth
cells over $z_{\max} = 1.0$~m.  The true profile is an exponential
with relaxation length 3~cm.

The scale factor $2 \times 10^4$ converts from Bq/m$^3$ to approximate
counts-per-second for a typical HPGe detector.

\subsection{Kernel Validation}

The numerical kernel (Gauss--Legendre quadrature with 4096 log-spaced
points over $u \in [1, 2\times 10^3]$) agrees with the analytic E1 kernel
to within 0.2\% relative error across all depths.  This validates the
discretisation and confirms that the Taylor buildup factor and angular
efficiency functions are correctly implemented as modular extensions.

\subsection{Non-Parametric Recovery}

For Scenario~A, Tikhonov/GCV recovers the areal activity within 30\%
relative error and the median depth $z_{\mathrm{med}}$ within 6~cm absolute
error, with $\chi^2/m < 2$.  This establishes the baseline performance
for the well-determined multi-line, multi-height configuration.

\subsection{Parametric Recovery}

For Scenario~A, the parametric pulse inversion recovers the deposition
density $A_0$ within 50\% and the effective time $t$ within 50\%.  The
wider tolerance reflects the strong correlation between $A_0$ and $t$
in the forward model: increasing $A_0$ and decreasing $t$ can produce
nearly identical profiles.  This is an inherent limitation of the
parametric approach that motivates the non-parametric and ensemble
methods.

\subsection{Solver Comparison}

All solvers (Tikhonov, Landweber, Cimmino, CGLS, Kaczmarz, FISTA, TV,
focusing MS/MGS) produce non-negative solutions with positive areal
activity.  The geophysical solvers exhibit characteristic differences:

\begin{itemize}
  \item \textbf{TV/ADMM} produces sharper layer boundaries than Tikhonov,
        at the cost of potential staircase artefacts in smooth regions.
        This is advantageous for scenarios with genuine discontinuities
        (buried layers, excavation boundaries).
  \item \textbf{FISTA} ($\ell_1$) produces the sparsest solutions, with
        activity concentrated in fewer depth cells than Tikhonov.
        The sparsity is controlled by the $\ell_1$ weight $\alpha$.
  \item \textbf{Focusing MGS} produces blocky, layered models with
        sharp interfaces and flat interiors, closely matching the
        expected geological structure of contaminated sites.
  \item \textbf{Focusing MS} produces compact solutions with isolated
        activity peaks, suitable for point-source contamination.
  \item \textbf{CGLS} converges in approximately 20--50 iterations
        for the typical problem size, with the L-curve corner
        providing a robust automatic stopping criterion.
\end{itemize}

\subsection{Criteria Comparison}

All six criteria (GCV, L-curve, discrepancy, quasi-optimality, NCP, SNR)
return finite, positive $\alpha$ values.  In practice, GCV and L-curve
provide the most robust selections for well-conditioned multi-line data,
while quasi-optimality and SNR are more reliable for poorly conditioned
single-line configurations where the noise level is uncertain.

\subsection{Transport Chemistry Validation}

The transport module is validated through internal consistency checks:
(1)~retardation factors agree with the analytical formula $R = 1 + \rho_b K_d/\theta$;
(2)~pH-dependent $K_d$ increases with pH for Cs-137 (consistent with
FES site deprotonation); (3)~Freundlich $K_d$ increases at lower
concentrations ($n_F < 1$); (4)~competitive $K_d$ decreases with
competitor ion concentration; (5)~U $K_d$ is higher under reducing
conditions; (6)~two-site retardation interpolates between the
instantaneous equilibrium value ($t \to 0$) and the full
retardation ($t \to \infty$); (7)~the Crank--Nicolson solver produces
non-negative concentrations in multi-layer configurations.

\subsection{Bayesian Uncertainty}

EKI with 50 ensemble members and 10 iterations produces non-negative
ensemble mean with finite per-depth uncertainty.  The posterior standard
deviation increases with depth, reflecting the decreasing sensitivity of
the kernel.  The Laplace MAP solution agrees with Tikhonov (as expected
from the Gaussian equivalence) but additionally provides the full
posterior covariance matrix for uncertainty propagation to derived
quantities (areal activity, median depth).


% =================================================================
\section{Discussion}\label{sec:discussion}

\subsection{Relation to the Geophysical Inversion Literature}

The adaptation of geophysical inversion methods to radionuclide depth
profiling is natural given the structural similarity of the problems.
Both involve linear (or linearised) forward operators with decaying
sensitivity, Poisson or Gaussian noise, and the need for
regularisation to obtain physically meaningful solutions.  However,
several aspects are specific to the gamma-spectrometry problem:

\begin{enumerate}
  \item \textbf{Strong exponential decay:} The E1 kernel decays as
  $e^{-\mu_s z}/z$ (asymptotically), which is faster than the
  algebraic decay ($1/z^3$) of gravity kernels.  This makes the
  ill-posedness more severe and amplifies the need for depth weighting.

  \item \textbf{Non-negativity constraint:} Activity cannot be negative,
  and this physical constraint significantly improves solution quality
  compared to unconstrained inversion.  All solvers in the framework
  enforce this by default.

  \item \textbf{Poisson noise:} Count data follow Poisson statistics
  ($\sigma_i = \sqrt{C_i}$), leading to heteroscedastic noise.  The
  weighted system formulation handles this naturally.

  \item \textbf{Transport-chemistry prior:} Unlike most geophysical
  inversions where the prior is purely mathematical (smoothness,
  compactness), the ADE provides a physics-based prior that directly
  encodes the migration behaviour of the contaminant.  The parametric
  inversion exploits this to extract physically meaningful parameters.
\end{enumerate}

\subsection{Choice of Solver for Different Scenarios}

The ensemble multi-method approach with AIC/BIC selection is recommended
for general use, as it automates the solver choice.  However, physical
considerations can guide the selection:

\begin{itemize}
  \item \textbf{Smooth, diffusive profiles} (e.g., Chernobyl Cs-137
  after decades of migration): Tikhonov with smoothness prior.
  \item \textbf{Layered contamination} (e.g., buried waste, plume
  interfaces): TV/ADMM or focusing MGS.
  \item \textbf{Compact, localised sources} (e.g., buried spent fuel
  fragments): FISTA ($\ell_1$) or focusing MS.
  \item \textbf{Uncertainty quantification}: EKI for full posterior
  bands; Laplace MAP for quick linearised estimates.
  \item \textbf{Well-understood transport history}: parametric pulse
  or exponential inversion for interpretable parameters.
  \item \textbf{Multiple nuclides}: joint inversion with coupling
  constraints to exploit shared depth structure.
\end{itemize}

\subsection{Limitations and Future Work}

Several limitations should be acknowledged.  First, the framework
assumes a laterally homogeneous activity distribution --- lateral
heterogeneity is not considered.  For sites with significant lateral
variability, a 2D or 3D inversion would be required, though the
computational cost would increase substantially.  Second, the
non-negativity constraint is implemented via simple projection
(clipping to zero), which can introduce bias near the boundary of
the feasible set; more sophisticated approaches (e.g., interior-point
methods) could improve this.  Third, the EKI uncertainty estimates
rely on the ensemble approximation and may underestimate posterior
variance for strongly non-Gaussian posteriors; comparison with MCMC
benchmarking would be valuable.  Fourth, the Crank--Nicolson solver is
limited to 1D; 2D migration modelling would require finite-element or
finite-volume methods.

Future directions include: (i)~integration of machine-learning
surrogates for fast forward modelling in 2D/3D; (ii)~adaptive mesh
refinement based on the resolution analysis; (iii)~sequential
assimilation of time-lapse measurements via the ensemble Kalman
filter; (iv)~extension to spectral unmixing of overlapping peaks
from multiple nuclides.

\subsection{Comparison with Existing Software}

The most widely used existing software for radionuclide depth profiling
is the IAEA's unreleased codes accompanying \citet{zombori1992new},
which implement basic unfolding methods (matrix inversion, strip
parameters) but lack modern regularisation.  \citet{tyler2008determination}
and \citet{hasan2022regularised,hasan2023bayesian} have developed
regularised and Bayesian approaches for borehole gamma spectrometry,
but their codes are not publicly available as a unified framework.  The
\texttt{soilactivity} package fills this gap by providing, for the first
time, a comprehensive, open-source Python framework that integrates
geophysical inversion methods with reactive transport chemistry for
in-situ gamma-spectrometry depth profiling.


% =================================================================
\section{Conclusions}\label{sec:conclusions}

This paper has presented the \texttt{soilactivity} software framework
for reconstructing vertical radionuclide activity profiles from in-situ
gamma-spectrometry measurements.  The key scientific contributions and
findings are:

\begin{enumerate}
  \item The systematic adaptation of ten geophysical inversion algorithms
  (Tikhonov, TSVD, Landweber, Cimmino, CGLS, Kaczmarz, FISTA, TV/ADMM,
  focusing MS/MGS) to the Fredholm gamma-transport equation
  demonstrates that methods developed for gravity, magnetic, and
  electromagnetic inversion transfer effectively to radionuclide
  depth profiling, with TV and focusing MGS providing superior
  recovery of piecewise-constant and layered contamination profiles.

  \item Six parameter selection criteria (GCV, L-curve, discrepancy,
  quasi-optimality, NCP, SNR) provide redundant and complementary
  approaches to the regularisation parameter choice, with GCV and
  L-curve being most robust for well-determined configurations and
  quasi-optimality being most robust for poorly conditioned ones.

  \item The integration of reactive transport chemistry --- including
  pH-dependent $K_d$, Freundlich isotherms, competitive ion exchange,
  Eh-dependent uranium speciation, two-site kinetic sorption, and
  multi-layer soil columns --- provides physics-based priors that
  yield directly interpretable inversion parameters and constrain
  the solution space beyond what purely mathematical regularisation
  can achieve.

  \item Bayesian uncertainty quantification via EKI delivers posterior
  uncertainty bands in seconds (versus hours for MCMC), enabling
  practical uncertainty propagation for field measurements.

  \item Depth weighting strategies from gravity/magnetic inversion
  (Li--Oldenburg, power-law, logarithmic, adaptive) effectively
  counteract the exponential sensitivity decay of the gamma kernel,
  improving the recovery of deep activity.

  \item Joint multi-nuclide inversion with coupling constraints and
  ensemble AIC/BIC model selection automates the inversion workflow
  and provides a principled approach to solver selection.
\end{enumerate}

The framework is validated through 25+ twin experiments covering all
solvers, criteria, transport models, diagnostics, and the full
pipeline.  The software is released as a Python package, designed for
practical deployment in environmental monitoring and decommissioning
programmes.


% =================================================================
\section*{Acknowledgements}

This work was supported by the \texttt{soilactivity} open-source
project.  The authors thank the developers of NumPy, SciPy, and
the broader Python scientific computing ecosystem.

% =================================================================
\bibliographystyle{apalike}
\bibliography{refs}

\end{document}
